\section{Invertible modules}
\label{section-invertible}

\noindent
Similarly to the case of modules over rings
(More on Algebra, Section \ref{more-algebra-section-picard})
we have the following definition.

\begin{definition}
\label{definition-invertible}
Let $(X, \mathcal{O}_X)$ be a ringed space. An
{\it invertible $\mathcal{O}_X$-module} is a sheaf
of $\mathcal{O}_X$-modules $\mathcal{L}$ such that
the functor
$$
\textit{Mod}(\mathcal{O}_X) \longrightarrow \textit{Mod}(\mathcal{O}_X),\quad
\mathcal{F} \longmapsto \mathcal{L} \otimes_{\mathcal{O}_X} \mathcal{F}
$$
is an equivalence of categories. We say that $\mathcal{L}$ is
{\it trivial} if it is isomorphic as an $\mathcal{O}_X$-module
to $\mathcal{O}_X$.
\end{definition}

\noindent
Lemma \ref{lemma-invertible-is-locally-free-rank-1}
below explains the relationship with locally free modules
of rank $1$.

\begin{lemma}
\label{lemma-invertible}
Let $(X, \mathcal{O}_X)$ be a ringed space. Let $\mathcal{L}$
be an $\mathcal{O}_X$-module. Equivalent are
\begin{enumerate}
\item $\mathcal{L}$ is invertible, and
\item there exists an $\mathcal{O}_X$-module $\mathcal{N}$
such that
$\mathcal{L} \otimes_{\mathcal{O}_X} \mathcal{N} \cong \mathcal{O}_X$.
\end{enumerate}
In this case $\mathcal{L}$ is locally a direct summand of a finite free
$\mathcal{O}_X$-module and the module $\mathcal{N}$ in (2) is isomorphic to
$\SheafHom_{\mathcal{O}_X}(\mathcal{L}, \mathcal{O}_X)$.
\end{lemma}

\begin{proof}
Assume (1). Then the functor $- \otimes_{\mathcal{O}_X} \mathcal{L}$
is essentially surjective, hence there exists an $\mathcal{O}_X$-module
$\mathcal{N}$ as in (2). If (2) holds, then the functor
$- \otimes_{\mathcal{O}_X} \mathcal{N}$ is a quasi-inverse
to the functor $- \otimes_{\mathcal{O}_X} \mathcal{L}$ and
we see that (1) holds.

\medskip\noindent
Assume (1) and (2) hold. Denote
$\psi : \mathcal{L} \otimes_{\mathcal{O}_X} \mathcal{N} \to \mathcal{O}_X$
the given isomorphism. Let $x \in X$. Choose an open neighbourhood
$U$ an integer $n \geq 1$ and sections $s_i \in \mathcal{L}(U)$,
$t_i \in \mathcal{N}(U)$ such that $\psi(\sum s_i \otimes t_i) = 1$.
Consider the isomorphisms
$$
\mathcal{L}|_U \to
\mathcal{L}|_U \otimes_{\mathcal{O}_U}
\mathcal{L}|_U \otimes_{\mathcal{O}_U} \mathcal{N}|_U \to \mathcal{L}|_U
$$
where the first arrow sends $s$ to $\sum s_i \otimes s \otimes t_i$
and the second arrow sends $s \otimes s' \otimes t$ to $\psi(s' \otimes t)s$.
We conclude that $s \mapsto \sum \psi(s \otimes t_i)s_i$ is
an automorphism of $\mathcal{L}|_U$. This automorphism factors as
$$
\mathcal{L}|_U \to \mathcal{O}_U^{\oplus n} \to \mathcal{L}|_U
$$
where the first arrow is given by
$s \mapsto (\psi(s \otimes t_1), \ldots, \psi(s \otimes t_n))$
and the second arrow by $(a_1, \ldots, a_n) \mapsto \sum a_i s_i$.
In this way we conclude that $\mathcal{L}|_U$ is a direct summand
of a finite free $\mathcal{O}_U$-module.

\medskip\noindent
Assume (1) and (2) hold. Consider the evaluation map
$$
\mathcal{L} \otimes_{\mathcal{O}_X}
\SheafHom_{\mathcal{O}_X}(\mathcal{L}, \mathcal{O}_X)
\longrightarrow \mathcal{O}_X
$$
To finish the proof of the lemma
we will show this is an isomorphism by checking it induces
isomorphisms on stalks. Let $x \in X$.
Since we know (by the previous paragraph)
that $\mathcal{L}$ is a finitely presented
$\mathcal{O}_X$-module
we can use Lemma \ref{lemma-stalk-internal-hom}
to see that it suffices to show that
$$
\mathcal{L}_x \otimes_{\mathcal{O}_{X, x}}
\Hom_{\mathcal{O}_{X, x}}(\mathcal{L}_x, \mathcal{O}_{X, x})
\longrightarrow \mathcal{O}_{X, x}
$$
is an isomorphism. Since
$\mathcal{L}_x \otimes_{\mathcal{O}_{X, x}} \mathcal{N}_x =
(\mathcal{L} \otimes_{\mathcal{O}_X} \mathcal{N})_x =
\mathcal{O}_{X, x}$ (Lemma \ref{lemma-stalk-tensor-product})
the desired result follows from
More on Algebra, Lemma \ref{more-algebra-lemma-invertible}.
\end{proof}

\begin{lemma}
\label{lemma-pullback-invertible}
Let $f : (X, \mathcal{O}_X) \to (Y, \mathcal{O}_Y)$ be a
morphism of ringed spaces. The pullback $f^*\mathcal{L}$ of an
invertible $\mathcal{O}_Y$-module is invertible.
\end{lemma}

\begin{proof}
By Lemma \ref{lemma-invertible}
there exists an $\mathcal{O}_Y$-module $\mathcal{N}$ such that
$\mathcal{L} \otimes_{\mathcal{O}_Y} \mathcal{N} \cong \mathcal{O}_Y$.
Pulling back we get
$f^*\mathcal{L} \otimes_{\mathcal{O}_X} f^*\mathcal{N} \cong \mathcal{O}_X$
by Lemma \ref{lemma-tensor-product-pullback}.
Thus $f^*\mathcal{L}$ is invertible by Lemma \ref{lemma-invertible}.
\end{proof}

\begin{lemma}
\label{lemma-invertible-is-locally-free-rank-1}
Let $(X, \mathcal{O}_X)$ be a ringed space. Any locally free
$\mathcal{O}_X$-module of rank $1$ is invertible.
If all stalks $\mathcal{O}_{X, x}$ are local rings, then
the converse holds as well (but in general this is not the case).
\end{lemma}

\begin{proof}
The parenthetical statement follows by considering a one point
space $X$ with sheaf of rings $\mathcal{O}_X$ given by a ring $R$.
Then invertible $\mathcal{O}_X$-modules correspond to invertible
$R$-modules, hence as soon as $\Pic(R)$ is not the trivial group,
then we get an example.

\medskip\noindent
Assume $\mathcal{L}$ is locally free of rank $1$ and consider the
evaluation map
$$
\mathcal{L} \otimes_{\mathcal{O}_X}
\SheafHom_{\mathcal{O}_X}(\mathcal{L}, \mathcal{O}_X)
\longrightarrow \mathcal{O}_X
$$
Looking over an open covering trivialization $\mathcal{L}$, we see
that this map is an isomorphism. Hence $\mathcal{L}$ is invertible
by Lemma \ref{lemma-invertible}.

\medskip\noindent
Assume all stalks $\mathcal{O}_{X, x}$ are local rings and $\mathcal{L}$
invertible. In the proof of Lemma \ref{lemma-invertible}
we have seen that $\mathcal{L}_x$ is an invertible
$\mathcal{O}_{X, x}$-module for all $x \in X$. Since
$\mathcal{O}_{X, x}$ is local, we see that
$\mathcal{L}_x \cong \mathcal{O}_{X, x}$
(More on Algebra, Section \ref{more-algebra-section-picard}).
Since $\mathcal{L}$ is of finite presentation by
Lemma \ref{lemma-invertible} we conclude that $\mathcal{L}$
is locally free of rank $1$ by
Lemma \ref{lemma-finite-presentation-stalk-free}.
\end{proof}

\begin{lemma}
\label{lemma-constructions-invertible}
Let $(X, \mathcal{O}_X)$ be a ringed space.
\begin{enumerate}
\item If $\mathcal{L}$, $\mathcal{N}$ are invertible
$\mathcal{O}_X$-modules, then so is
$\mathcal{L} \otimes_{\mathcal{O}_X} \mathcal{N}$.
\item If $\mathcal{L}$ is an invertible $\mathcal{O}_X$-module, then so is
$\SheafHom_{\mathcal{O}_X}(\mathcal{L}, \mathcal{O}_X)$ and the evaluation map
$\mathcal{L} \otimes_{\mathcal{O}_X}
\SheafHom_{\mathcal{O}_X}(\mathcal{L}, \mathcal{O}_X) \to \mathcal{O}_X$
is an isomorphism.
\end{enumerate}
\end{lemma}

\begin{proof}
Part (1) is clear from the definition and part (2) follows from
Lemma \ref{lemma-invertible} and its proof.
\end{proof}

\begin{definition}
\label{definition-powers}
Let $(X, \mathcal{O}_X)$ be a ringed space. Given an invertible sheaf
$\mathcal{L}$ on $X$ and $n \in \mathbf{Z}$ we define the
$n$th {\it tensor power} $\mathcal{L}^{\otimes n}$ of $\mathcal{L}$
as the image of $\mathcal{O}_X$ under applying the equivalence
$\mathcal{F} \mapsto\mathcal{F} \otimes_{\mathcal{O}_X} \mathcal{L}$
exactly $n$ times.
\end{definition}

\noindent
This makes sense also for negative $n$ as we've defined an invertible
$\mathcal{O}_X$-module as one for which tensoring is an equivalence.
More explicitly, we have
$$
\mathcal{L}^{\otimes n} =
\left\{
\begin{matrix}
\mathcal{O}_X & \text{if} & n = 0 \\
\SheafHom_{\mathcal{O}_X}(\mathcal{L}, \mathcal{O}_X) & \text{if} & n = -1\\
\mathcal{L} \otimes_{\mathcal{O}_X} \ldots \otimes_{\mathcal{O}_X} \mathcal{L}
& \text{if} & n > 0 \\
\mathcal{L}^{\otimes -1} \otimes_{\mathcal{O}_X} \ldots
\otimes_{\mathcal{O}_X} \mathcal{L}^{\otimes -1}
& \text{if} & n < -1
\end{matrix}
\right.
$$
see Lemma \ref{lemma-constructions-invertible}.
With this definition we have canonical isomorphisms
$\mathcal{L}^{\otimes n} \otimes_{\mathcal{O}_X}
\mathcal{L}^{\otimes m} \to
\mathcal{L}^{\otimes n + m}$, and these isomorphisms
satisfy a commutativity and an associativity constraint
(formulation omitted).

\medskip\noindent
Let $(X, \mathcal{O}_X)$ be a ringed space. We can define a $\mathbf{Z}$-graded
ring structure on $\bigoplus \Gamma(X, \mathcal{L}^{\otimes n})$ by mapping
$s \in \Gamma(X, \mathcal{L}^{\otimes n})$ and
$t \in \Gamma(X, \mathcal{L}^{\otimes m})$ to the
section corresponding to $s \otimes t$ in
$\Gamma(X, \mathcal{L}^{\otimes n + m})$.
We omit the verification that this defines a commutative
and associative ring with $1$. However, by our conventions in
Algebra, Section \ref{algebra-section-graded}
a graded ring has no nonzero elements in negative degrees.
This leads to the following definition.

\begin{definition}
\label{definition-gamma-star}
Let $(X, \mathcal{O}_X)$ be a ringed space.
Given an invertible sheaf $\mathcal{L}$ on $X$ we define
the {\it associated graded ring} to be
$$
\Gamma_*(X, \mathcal{L})
=
\bigoplus\nolimits_{n \geq 0} \Gamma(X, \mathcal{L}^{\otimes n})
$$
Given a sheaf of $\mathcal{O}_X$-modules $\mathcal{F}$ we set
$$
\Gamma_*(X, \mathcal{L}, \mathcal{F})
=
\bigoplus\nolimits_{n \in \mathbf{Z}} \Gamma(X,
\mathcal{F} \otimes_{\mathcal{O}_X} \mathcal{L}^{\otimes n})
$$
which we think of as a graded $\Gamma_*(X, \mathcal{L})$-module.
\end{definition}

\noindent
We often write simply $\Gamma_*(\mathcal{L})$ and $\Gamma_*(\mathcal{F})$
(although this is ambiguous if $\mathcal{F}$ is invertible).
The multiplication of $\Gamma_*(\mathcal{L})$ on
$\Gamma_*(\mathcal{F})$ is defined using the isomorphisms
above. If $\gamma : \mathcal{F} \to \mathcal{G}$ is a $\mathcal{O}_X$-module
map, then we get an $\Gamma_*(\mathcal{L})$-module homomorphism
$\gamma : \Gamma_*(\mathcal{F}) \to \Gamma_*(\mathcal{G})$.
If $\alpha : \mathcal{L} \to \mathcal{N}$ is an $\mathcal{O}_X$-module
map between invertible $\mathcal{O}_X$-modules, then we obtain
a graded ring homomorphism $\Gamma_*(\mathcal{L}) \to \Gamma_*(\mathcal{N})$.
If $f : (Y, \mathcal{O}_Y) \to (X, \mathcal{O}_X)$
is a morphism of ringed spaces and if $\mathcal{L}$ is invertible
on $X$, then we get an invertible sheaf $f^*\mathcal{L}$ on $Y$
(Lemma \ref{lemma-pullback-invertible})
and an induced homomorphism of graded rings
$$
f^* :
\Gamma_*(X, \mathcal{L})
\longrightarrow
\Gamma_*(Y, f^*\mathcal{L})
$$
Furthermore, there are some compatibilities between the constructions
above whose statements we omit.

\begin{lemma}
\label{lemma-pic-set}
Let $(X, \mathcal{O}_X)$ be a ringed space.
There exists a set of invertible modules $\{\mathcal{L}_i\}_{i \in I}$
such that each invertible module on $X$ is isomorphic to exactly
one of the $\mathcal{L}_i$.
\end{lemma}

\begin{proof}
Recall that any invertible $\mathcal{O}_X$-module is locally
a direct summand of a finite free $\mathcal{O}_X$-module, see
Lemma \ref{lemma-invertible}.
For each open covering $\mathcal{U} : X = \bigcup_{j \in J} U_j$
and map $r : J \to \mathbf{N}$ consider the sheaves of
$\mathcal{O}_X$-modules $\mathcal{F}$ such that
$\mathcal{F}_j = \mathcal{F}|_{U_j}$ is a direct summand of
$\mathcal{O}_{U_j}^{\oplus r(j)}$.
The collection of isomorphism classes of $\mathcal{F}_j$ is a set, because
$\Hom_{\mathcal{O}_U}(\mathcal{O}_U^{\oplus r}, \mathcal{O}_U^{\oplus r})$
is a set. The sheaf $\mathcal{F}$ is gotten by glueing $\mathcal{F}_j$,
see Sheaves, Section
\ref{sheaves-section-glueing-sheaves}. Note that the collection of
all glueing data forms a set. The collection of all coverings
$\mathcal{U} : X = \bigcup_{j \in J} U_i$ where $J \to \mathcal{P}(X)$,
$j \mapsto U_j$ is injective forms a set as well. For each covering
there is a set of maps $r : J \to \mathbf{N}$. Hence the collection
of all $\mathcal{F}$ forms a set.
\end{proof}

\noindent
This lemma says roughly speaking that the collection of
isomorphism classes of invertible sheaves forms a set.
Lemma \ref{lemma-constructions-invertible} says that
tensor product defines the structure of an abelian group
on this set.

\begin{definition}
\label{definition-pic}
Let $(X, \mathcal{O}_X)$ be a ringed space.
The {\it Picard group} $\Pic(X)$ of $X$ is the
abelian group whose elements are isomorphism classes of
invertible $\mathcal{O}_X$-modules, with addition
corresponding to tensor product.
\end{definition}

\begin{lemma}
\label{lemma-s-open}
\begin{slogan}
A (local) trivialisation of a linebundle
is the same as a (local) nonvanishing section.
\end{slogan}
Let $X$ be a ringed space. Assume that each stalk $\mathcal{O}_{X, x}$
is a local ring with maximal ideal $\mathfrak m_x$.
Let $\mathcal{L}$ be an invertible $\mathcal{O}_X$-module.
For any section $s \in \Gamma(X, \mathcal{L})$ the set
$$
X_s = \{x \in X \mid \text{image }s \not\in \mathfrak m_x\mathcal{L}_x\}
$$
is open in $X$. The map $s : \mathcal{O}_{X_s} \to \mathcal{L}|_{X_s}$
is an isomorphism, and there exists a section $s'$
of $\mathcal{L}^{\otimes -1}$ over $X_s$ such that $s' (s|_{X_s}) = 1$.
\end{lemma}

\begin{proof}
Suppose $x \in X_s$.
We have an isomorphism
$$
\mathcal{L}_x \otimes_{\mathcal{O}_{X, x}} (\mathcal{L}^{\otimes -1})_x
\longrightarrow
\mathcal{O}_{X, x}
$$
by Lemma \ref{lemma-constructions-invertible}.
Both $\mathcal{L}_x$ and $(\mathcal{L}^{\otimes -1})_x$
are free $\mathcal{O}_{X, x}$-modules of rank $1$. We conclude
from Algebra, Nakayama's Lemma \ref{algebra-lemma-NAK} that
$s_x$ is a basis for $\mathcal{L}_x$. Hence there exists
a basis element $t_x \in (\mathcal{L}^{\otimes -1})_x$
such that $s_x \otimes t_x$ maps to $1$.
Choose an open neighbourhood $U$ of
$x$ such that $t_x$ comes from a section $t$
of $\mathcal{L}^{\otimes -1}$ over $U$ and such that
$s \otimes t$ maps to $1 \in \mathcal{O}_X(U)$.
Clearly, for every $x' \in U$ we see that $s$ generates
the module $\mathcal{L}_{x'}$. Hence $U \subset X_s$.
This proves that $X_s$ is open. Moreover, the section
$t$ constructed over $U$ above is unique, and hence
these glue to give the section $s'$ of the lemma.
\end{proof}

\begin{remark}
\label{remark-pullback-s-open}
Given a morphism of locally ringed spaces $f : Y \to X$
(see Schemes, Definition \ref{schemes-definition-locally-ringed-space})
the inverse image $f^{-1}(X_s)$ is equal to $Y_{f^*s}$, where
$f^*s \in \Gamma(Y, f^*\mathcal{L})$ is the pullback of $s$.
\end{remark}


